\documentclass[12pt, b5paper]{article}
\setlength{\topmargin}{-.5in}
\setlength{\textwidth}{6.5in}
\setlength{\textheight}{9.5in}
\setlength{\oddsidemargin}{-.25in}

\usepackage{amsmath}

\newcommand{\floor}[1]{\left\lfloor #1 \right\rfloor}

\pagestyle{empty}
\begin{document}

\begin{center}
{\bf CSC 226
Exercise 3}\\
An Nguyen, V01073373\\
22 Feb 2026
\end{center}

\section*{Problem:}
Prove the theorem: Let G be a connected edge-weighted undirected graph with distinct edge weights, 
and let C be some cycle in G. Then the maximum-weight edge e in C is not contained in the 
minimum weight spanning tree of G.

\section*{Proof:}
We will prove the theorem by contradiction. We know that in a connected graph G, where all edge weights are distinct, there is 
a unique minimum weight spanning tree (MWST). If an edge satisfies to be the minimum 
weight bridge for a specific cut, it must be in that unique tree.
\subsubsection*{Claim:} 
Let C be a cycle in G, and let e be the maximum-weight edge in that cycle. 
We want to prove that e is not in the MWST.
\subsubsection*{Assumption:} 
Suppose, for the purpose of contradiction, e is in the MWST. Deleting any edge from a spanning tree disconnects it into exactly 
two connected components, which we can call S and V - S. By deleting e from our assumed tree, 
we create such a partition, or a cut, in the graph.
Since e was in the cycle C, there must be another edge f in C that connects S and V - S so that the
cycle can cross the cut at e and go through the cut again at f.\\
\newline
From the assumption, we have that e is the maximum-weight edge in C. Since all edge weights in G
 are distinct, any other edge in C, including f, must have a weight less than that of e.\\
\newline
By the Cut Property, for any partition of a graph, the minimum weight edge crossing 
the cut must be in the MWST. In our case S and V - S form a cut, and we have edge f 
crossing the cut with weight less than that of e. Therefore, edge f must be in the MWST. Since 
we also assumed that e is in the MWST, we have a contradiction: a spanning tree cannot contain a cycle.
\subsubsection*{Conclusion:}
Thus, our assumption that e is in the MWST must be false. Therefore, the maximum-weight edge e 
in cycle C is not contained in the minimum weight spanning tree of G.\\
\end{document}